\usepackage[utf8]{inputenc}
\usepackage[T1]{fontenc} % optional
\usepackage{amsmath,amsfonts}
\usepackage{bm}
\usepackage[scale=.9]{cascadia-code}
\usepackage{algorithmic}
\usepackage{graphicx}
\usepackage{textcomp}
\usepackage{xfp}
\usepackage{xcolor}
\usepackage{colortbl}
\usepackage[backend=biber,style=acmnumeric]{biblatex}
\addbibresource{references.bib}
\AtBeginDocument{%
  \providecommand\BibTeX{{%
    Bib\TeX}}}
\usepackage{url}
\definecolor{spaceorange}{HTML}{ED820E}
\long\def\todo#1{{\tiny\sffamily\color{purple!75!black}{\bfseries\faCaretRight\space ToDo:}~#1~\textcolor{purple!75!black}{\faCaretLeft}}}
\def\plannedSpace#1{\vphantom{I}\rlap{\smash{\raisebox{1.66\ht\strutbox}{\parbox\linewidth{\strut\hfill\sffamily\color{spaceorange}#1\hfill\strut}}}}}
\usepackage{fontawesome}
\usepackage{tikz}
\usetikzlibrary{tikzmark,shapes.symbols}
\usepackage{makecell}
\usepackage{booktabs}
\usepackage{csquotes}
\PassOptionsToPackage{hidelinks}{hyperref}
\usepackage{hyperref}
\usepackage[capitalize,nameinlink,noabbrev]{cleveref}
\let\T\texttt
\let\say\enquote
\def\basecomment#1#2#3{\space\faCaretRight\;\textcolor{#1}{\textbf{#2:}~#3}\;\faCaretLeft\space}
\def\fs#1{\basecomment{teal}{fs}{#1}}
\def\og#1{\basecomment{green!60!purple!60!black}{og}{#1}}
\def\lp#1{\basecomment{cyan!60!purple!60!black}{lp}{#1}}
\def\mtt#1{\basecomment{purple!50!orange!60!black}{mtt}{#1}}
\usepackage{microtype}
\usepackage[inline]{enumitem}
\newlist{inlist}{enumerate*}{1}
\setlist[inlist]{itemjoin={{, }},itemjoin*={{, and }},label=(\(\roman*\)),mode=boxed}
\usepackage{forest}
\useforestlibrary{edges}
\def\SWidth{3.25cm}%
\forestset{
   internal/.style={gray,font=\sffamily},
   marker/.style={rectangle split parts=1,minimum width=0cm},
   norm-ast-node/.style={
      Rect,
      outer sep=1pt,
      rounded corners=1pt,
      Soft,
      s sep=2pt,
      l sep+=2mm,
      rectangle split,
      rectangle split parts=2,
      font=\ttfamily,
   },
   l2r/.style={
      grow'=east,
      anchor = west,
      parent anchor = east,
      child anchor = west,
   },
   rot90/.style={
      rotate=90,grow=east,anchor=north,parent anchor=south,child anchor=north
   },
   norm-ast/.style={
      norm-ast-node,
      minimum width=\SWidth,
      l2r,
      line join=round,
   },
   o/.style={edge label={node[pos=.5,fill=white,inner sep=1.25pt,rounded corners=1pt,font=\scriptsize] {\strut #1}}},
   al-center/.style={
      before computing xy={s/.average={s}{siblings}}
   }
}
\def\type#1{:~{\textit{\smaller[2]#1}}}%
\def\NodeSurround{~~}%
\newsavebox\NodeTitleBox
\newcommand\Node[3][\SWidth]{\global\setbox\NodeTitleBox=\hbox{\NodeSurround#2\NodeSurround}\usebox\NodeTitleBox \nodepart{two} \edef\w{\ifdim\wd\NodeTitleBox<#1 #1\else\wd\NodeTitleBox\fi}\parbox{\w}{\smaller\sffamily#3}}%
\def\Or{\,|\,}
\newlist{blocklist}{enumerate*}{1}
\setlist[blocklist]{itemjoin={{\space}},itemjoin*={{\space}},label=(\(\roman*\)),mode=boxed}
\newlist{orlist}{enumerate*}{1}
\setlist[orlist]{itemjoin={{, }},itemjoin*={{, or }},label=(\(\alph*\)),mode=boxed}
\usepackage{wrapfig}
\usetikzlibrary{arrows.meta,fit,backgrounds,calc}
\tikzset{
   Soft/.style={line join=round,line cap=round},
   All Soft/.style={every path/.append style={Soft}},
   Blob/.style={
      draw=black,
      circle,
      outer sep=2pt,
      minimum size=1.8em
   },
   Blobs/.style={
      every node/.append style={Blob}
   },
   Use/.style={Blob},
   RectRounding/.style={rounded corners=3pt},
   Rect/.style={
      draw=black,
      rectangle,
      RectRounding,
      outer sep=2pt,
      minimum size=1.8em
   },
   Def/.style={Rect},
   Rects/.style={
      every node/.append style={Rect,minimum width=#1}
   },
   Rects/.default={1.5em},
   Link/.style={
      draw,
      Soft,
      rounded corners=2pt,
      -Kite
   },
   Links/.style={
      every path/.append style={Link}
   },
   comm/.style={rectangle,draw,text width=7.5mm,align=center,minimum height=5mm,font=\small\ttfamily,fill=lightgray!22!white},
   d/.style={comm,rounded corners=2pt,hyperlink node={v:def}}, % variable def
   u/.style={comm,rounded corners=2.5mm,hyperlink node={v:use}}, % variable use % rounded rectangle breaks anchors
   v/.style={comm,signal, signal to=east and west, gray, text=black, inner sep=-4pt, text width=5mm,fill=lightgray!22!white,hyperlink node={v:value}}, % values
   F/.style={comm,draw=gray,rounded corners=2pt,fill=white,inner xsep=.5em,hyperlink node={v:fdef}},
   fc/.style={comm,double,rounded corners=2.5mm, outer sep=1.15pt,hyperlink node={v:call}},
   w-back/.style={fill=white,inner sep=1pt},
   T/.style={font=\footnotesize\ttfamily,text=gray},
   code/.style={font=\ttfamily},
   dfidn/.style={
      circle,darkgray,fill opacity=.925,xshift=-1mm,yshift=.15mm,fill=white,draw,minimum size=1.8em,scale=.8,inner sep=-1pt
   },
   klabel/.style={font=\scriptsize\sffamily, inner sep=1pt,text=black,above,execute at begin node={\strut}},
   olabel/.style={midway,klabel},
   pos at/.style={above #1=-.5mm,yshift=-.4\baselineskip},
   old/.style={opacity=.65},
   graph-frame-style/.style={very thick,rounded corners=1.5pt,draw,black},
   graph-cd/.style={gray,font=\scriptsize,line cap=round}, 
   graph-cd-in/.style={graph-cd,{Kite[scale=.7]}-},
   graph-cd-out/.style={graph-cd,-{Kite[scale=.7]}}
}
\tikzset{
    hyperlink node/.style={
        alias=sourcenode,
        append after command={
            let     \p1 = (sourcenode.north west),
                \p2=(sourcenode.south east),
                \n1={\x2-\x1},
                \n2={\y1-\y2} in
            node [inner sep=0pt, outer sep=0pt,anchor=north west,at=(\p1)] {\hyperlink{#1}{\XeTeXLinkBox{\phantom{\rule{\n1}{\n2}}}}}
        }
    }
}
\def\DefineEdgeType#1#2#3{%
   \expandafter\def\csname edge#1short\endcsname{#1}
   \expandafter\def\csname edge#1long\endcsname{#2}
   \expandafter\def\csname edge#1desc\endcsname{\def\s{\textit{source}\xspace}\def\t{\textit{target}\xspace}#3}
}
\usepackage{xspace}
\def\GetEdge#1#2{\csname edge#1#2\endcsname}
\def\edge#1{\hyperlink{e:#1}{\textsf{\GetEdge{#1}{short}}}}
\def\E#1{\textit{\edge{#1}}}
\def\V#1{\hyperlink{v:#1}{\textsf{\textit{#1}}}}
\DefineEdgeType{reads}{reads}{\s reads from \t}
\DefineEdgeType{def-by}{defined-by}{\s is defined by the \t}
\DefineEdgeType{calls}{calls}{\s calls the \t}
\DefineEdgeType{returns}{returns}{\s returns to \t}
\DefineEdgeType{def-on-call}{defines (on call)}{if the call occurs, \s~(argument) defines \t~(parameter)} % . This edge is bi-directional.
\DefineEdgeType{arg}{argument}{\t vertex is an argument of the \s}
\DefineEdgeType{side-effect}{side effect (on-call)}{if the call occurs, the \s effects on \t}
\DefineEdgeType{nse}{non-standard evaluation}{\s causes a non-standard evaluation of \t}
\RequirePackage[
   print,
   numinpar,
   fakeminted
]{lib/xlistings/xlistings}
\lstcolorlet{keywordA}{black}
\lstcolorlet{keywordB}{black}
\lstcolorlet{keywordC}{black}
\lstcolorlet{literals}{black}
\lstcolorlet{numbers}{darkgray}
\xlstDefineStyles{%
 {keywordA: \color{xlst@colors@lst@keywordA}\bfseries},%
 {keywordB: \color{xlst@colors@lst@keywordB}},%
 {keywordC: \color{xlst@colors@lst@keywordC}},%
 {keywordD: \itshape},%
 {numbers:  \color{xlst@colors@lst@numbers}\slshape},%
 {literals: \color{xlst@colors@lst@literals}\slshape},%
 {comments: \color{xlst@colors@lst@comments}},% \scshape
 {basic:   \ttfamily},%
 {command:  \color{xlst@colors@lst@command}}%
}
\LoadLanguage{sle-R}
\xlstsetmintedstyle{plain number}
\def\LeftArrow{\text{\BeginAccSupp{method=escape,ActualText={<-}}\(\leftarrow\)\EndAccSupp{}}}
\def\RightArrow{\text{\BeginAccSupp{method=escape,ActualText={->}}\(\rightarrow\)\EndAccSupp{}}}
\def\DoubleLeftArrow{\text{\BeginAccSupp{method=escape,ActualText={<<-}}\(\twoheadleftarrow\)\EndAccSupp{}}}
\def\DoubleRightArrow{\text{\BeginAccSupp{method=escape,ActualText={->>}}\(\twoheadrightarrow\)\EndAccSupp{}}}
\AtBeginDocument{\xlstInlinePreBreak{}\xlstPreBreak{}}
\RequirePackage{tabularx}
\RequirePackage[detect-all]{siunitx}
\sisetup{number-mode=match,round-mode=places,round-pad=false,group-minimum-digits=4,group-separator={,}}%
\usepackage{multicol}
\usepackage{xfp}
\usepackage{multirow}
\def\PrintResultOfThe#1#2#3#4{%
   \edef\result{\fpeval{round(#1,#2)}}%
   \edef\chk{\fpeval{\result=0?1:0}}%
   \ifnum\chk=1\relax
      \edef\numcmpres{\fpeval{1/(10^#2)}}%
      \qty{< \numcmpres}{#4}%
   \else
      #3\relax
   \fi
}
\newcommand*\ThePercent[3][2]{%
   \edef\fst{\fpeval{#2}}%
   \edef\snd{\fpeval{#3}}%
   \edef\result{\fpeval{\fst/\snd * 100}}%
   \PrintResultOfThe{\result}{#1}{\ifdim\result sp=100sp\relax\ifdim\fst sp=\snd sp\else\(\approx\)\fi\fi\qty{\fpeval{\result}}{\percent}}{\percent}%
}
\def\FAIR{\textsc{fair}\xspace}
\usepackage[subtle]{savetrees}
